% Standalone increment: no prior source is modified.
\documentclass[11pt]{article}
\usepackage[T1]{fontenc}\usepackage[utf8]{inputenc}\usepackage{lmodern}
\usepackage[margin=.86in]{geometry}
\usepackage{amsmath,amssymb,amsthm,mathtools,microtype,array,longtable,booktabs,xurl}
\usepackage[hidelinks]{hyperref}
\setlength{\emergencystretch}{3em}\allowdisplaybreaks[2]
\numberwithin{equation}{section}
\newtheorem{theorem}{Theorem}[section]
\newtheorem{proposition}[theorem]{Proposition}
\newtheorem{lemma}[theorem]{Lemma}
\newtheorem{corollary}[theorem]{Corollary}
\theoremstyle{remark}\newtheorem{remark}[theorem]{Remark}
\newcommand{\E}{\mathbb E}\newcommand{\dd}{\,d}
\newcommand{\one}{\mathbf1}\newcommand{\eps}{\varepsilon}
\newcommand{\KK}{\mathcal K}\newcommand{\BB}{\mathcal B}
\newcommand{\Gp}{{G_+}}\newcommand{\CC}{\mathcal C}
\newcommand{\Kfun}{\mathsf K}\newcommand{\Ffun}{\mathsf F}
\newcommand{\dist}{\mathfrak d}
\title{Compact-window interface and probe-precision audit\\
with a first-moment leading-population extension}
\author{Standalone analytical increment for independent review}
\date{Version 1 --- 2 November 2026}
\begin{document}\maketitle
\begin{abstract}
The selected entry at a fixed small logistic level supplies a
mass-weighted first-moment comparison, but not compact weak support
or an original-flow comparison on the learned windows. We identify the
correct normalized centering and classify probe activity, including
boundary labels and the distinction between ancestry and current angle.
At entry, the unpaid outside probe quantity is a mass-weighted squared
active input, at precision \(o(q^2)\); its source-rate analogue is
controlled by the same quantity under the passage bounds. Existing
mass and tail estimates do not discharge every such charge.
In the defined leading system, WIT's rate robustness and finite-window
propagation extend to finite first weak moments, with its original
constants. This yields a strong-only compact-support version of its
population witness implication. No original trajectory is continued
beyond its proved interval.
\end{abstract}

\section{Scope, conventions and the consumer inventory}
\label{an07cwiv1:scope}
Throughout \(\kappa=10\), \(\rho\in[13/20,7/10]\),
\(\lambda=\rho^2\), \(Q=\lambda+q^2\), \(\ell=\log(1/q)\),
and \(\chi=1/20\). Initial radial scale is \(S_0=q^{10}\).
The original balanced representation is
\[
 U=\sqrt m\,a_\theta,\quad W=\sqrt m\,a_\psi,\qquad
 a_\theta=(\cos\theta,\sin\theta),\qquad
 f_T(X)=\int_Tm a_\psi(a_\theta\cdot X)_+\dd\lambda_0 .
\]
All cohorts used in output or rate decompositions are fixed labels.
Retain the partition \(\KK,\Gp,\Sigma,\BB\), selected clock
\(S=\mathcal C_q\), \(P_q=\Gp\setminus S\), and the outside
decomposition of RB. In particular \(\CC\subset\BB\) denotes the
fixed strong complement, not all nonselected labels.
In this note \(M=\mu_\KK\) and \(N\) is a leading weak mass;
neither is the residual constant \(M_{\rm res}\).
The leading coordinates are
\[
 a_s=(\cos qx,\sin qx),\quad b_s=(\cos qy,\sin qy),\qquad
 a_w=(\sin qu,\cos qu),\quad b_w=(-\sin qv,\cos qv).
\]
Thus the opposite conventions for weak input and output are preserved.

Sources are cited by exact label suffixes: WIT has prefix
\texttt{inc:AN06:v2:}; LP \texttt{inc:LP:v1:}; LPR
\texttt{inc:LPR:v1:}; TAILREP \texttt{inc:AN03:tailrepair:v1:}.
WCS, RB, SRR, CFA and S2CP have the prefixes specified in the
roadmap's source ledger. H2 means the handoff, Section 3.2 and AN7.
The companion review records repository paths and fetched blob hashes.

\paragraph{Entry already proved, not assumed anew.}
Keep WCS's explicit level
\[
 n=\frac12\min\left\{h,\frac14,
 \frac{\lambda_{\min}g_{\min}}{64(980)B_{\max}(20)}\right\},
 \quad \lambda_{\min}=\frac{169}{400},\
 g_{\min}=\frac{1183}{41600},\
 B_{\max}=\mathrm e\,2^{1/(2\lambda_{\min})}.
\]
WCS \texttt{main,handoff} gives, at its exact auxiliary-logistic hit
\(t_n<t_e<t_b\),
\begin{align}
 E_s+W_2&=O_n(q^\chi),&
 r_\KK&=O(q^\chi),&
 M&=1+O(q^{11/10}), \notag\\
 N_S/n&=1+O_n(q^{1/8}),&
 c&=n+O_n(q^{1/8}),&
 \sup\mu_\BB/q&\le q^{1/10}.
 \label{an07cwiv1:entry}
\end{align}
The \(W_2\) law is the fixed incoming probability \(\eta_s\);
current probability differs by a uniformly \(1+o(1)\) density.
Consequently, with the physical current core and selected measures,
\begin{equation}
 \begin{split}
 \dist(t_n):={}&|M-1|+|N_S-n|\\
 &+\int_\KK(|x-X_*(n)|+|y-Y_*(n)|)\dd\mu\\
 &+\int_S(|u-U_*(n)|+|v-V_*(n)|)\dd\mu
 =O_n(q^\chi).
 \end{split}
 \label{an07cwiv1:entrydistance}
\end{equation}
This uses \(L^2(\eta_s)\) to \(L^1(\eta_s)\), not a support
estimate. WBR/WCS only give a growing weak chart
\(|u|+|v|\le Cq^{-p_w}\), with \(9/10<p_w<1\).

\subsection{Time and parameter range}
\begin{proposition}[An admissible comparison clock, not an unproved hit]
\label{an07cwiv1:time}
The times denoted by \(t\) in the passage sources obey
\(\tau=\mu_1^2 t_{\rm phys}\) (\(\tau:\text{ normalized time}\)). Set
\begin{equation}
 \tau_h=t_n+\lambda^{-1}\log\frac{1-n}{n},\qquad
 s=\tau-\tau_h,\qquad s_n=\lambda^{-1}\log\frac n{1-n}.
 \label{an07cwiv1:center}
\end{equation}
This is exactly the half-mass center of WCS's own auxiliary logistic
mass, at the same \(\rho\). It is analytically defined by the proved
hit \(t_n\). In LP use \(t_h^{\rm phys}=\tau_h/\mu_1^2\), so that
\[
 t_h^{\rm phys}+s/\mu_1^2=(\tau_h+s)/\mu_1^2,\qquad
 D_p^{\rm phys}=\mu_1^2D_p^\tau .
 \]
The probe remains a unit vector. Moreover
\begin{equation}
 \tau_h=t_0-Q^{-1}\log H_0+o(1),\qquad
 (Q^{-1}-\lambda^{-1})O(\ell)=O(q^2\ell)=o(1).
 \label{an07cwiv1:Qtime}
\end{equation}
\end{proposition}
\begin{proof}
WCS gives \(t_n-t_0=Q^{-1}\log\{n/[(1-n)H_0]\}+o(1)\).
Insert this in \eqref{an07cwiv1:center}; the remaining fixed-level
logarithm has coefficient \(Q^{-1}-\lambda^{-1}\). The exact
logistic law gives WIT \texttt{clockalignment}. Since
\(t_0=10\ell+O(1)\), \(H_0\asymp A_q^{-\lambda}q^{10(1-\lambda)}\)
and \(A_q\) is bounded above and below, \(\tau_h\to\infty\).
Every fixed centered interval therefore has nonnegative original
time for small \(q\).
\end{proof}
The physical mass \(N_S\) is not asserted logistic. A leading
comparison initialized with that mass has own center
\[
 t_n+\lambda^{-1}\log\frac{1-N_S(t_n)}{N_S(t_n)}
       =\tau_h+O_n(q^{1/8}).
 \]
Alternatively multiply its initial weak measure by \(n/N_S(t_n)\);
this constructs a leading comparison with center exactly \(\tau_h\)
and changes \eqref{an07cwiv1:entrydistance} by \(O_n(q^{1/8})\).
Neither construction proves an original half-mass hit.
For a perturbed mass, WIT's logit identity retains
\(\int r_N/[N(1-N)]\); centering cannot remove this defect.

WIT's common windows and LPR's combined theorem hold only on their
proved rectangle
\(\mathcal P_*\subset[13/20,7/10]\times[87/50,7/4]\), centered at
\((2/3,5\pi/9)\). They are not established on the entire ratio box
with common margins. On those windows
\(n_\rho(s)\ge1/4+\sigma/2>n\), so the windows are a fixed time
after \(s_n\). This is compatible with the centering; it does not
put them inside the weak passage. No comparison beyond \(t_b\)
is supplied by RB. A fixed interval beginning before \(s_n\)
requires matching from a smaller fixed level, or a separate
earlier-window transfer. Its auxiliary logistic center would
be the same by the exact logistic identity. The necessary
earlier-window matching is not asserted by this audit.

\subsection{Complete hypothesis inventory}
\label{an07cwiv1:inventory}
Here E means supplied at \(t_n\), P means proved only on a
passage subinterval (never beyond \(t_b\)), L means an existing
leading-model conclusion, and O means open for the initialized
application. Model hypotheses and original-to-model identification
are separate obligations.
\small
\begin{longtable}{>{\raggedright\arraybackslash}p{.19\textwidth}>{\raggedright\arraybackslash}p{.42\textwidth}>{\raggedright\arraybackslash}p{.31\textwidth}}
\toprule Consumer & Every required domain/input & Status and source\\\midrule
\endhead
WIT \texttt{ratebound}
& \(R\ge2,\ 0<\rho\le1\), nonnegative measures,
 \(M,N\le2\), coherent masses \(1,n_*\), \(0\le n_*\le1\);
 same \(\rho,\zeta\).
& E for \(\KK,S\) by WCS; L for the comparator.
 Not an exact original leading-system identification.\\
& Both supports and comparator coordinates in \([-R,R]^2\);
 \(|\zeta|\le R\); strict gap \(\zeta-X\ge g_0>0\).
& E for strong support and bounded atom coordinates; O for fixed
 weak support. Theorem \ref{an07cwiv1:rate} replaces just that
 hypothesis. Gap supplied at entry on WIT offsets and on \(J\)
 by WIT \texttt{uniformtheorem}, not for arbitrary offsets.\\
& Selector convention at \(x=\zeta\); same definitions of
 \(\mathcal E,d_1,d_2,\dist\), including mass and both first moments.
& L definitions; E distance \eqref{an07cwiv1:entrydistance}.
 Core self-error has these leading definitions; full original
 observable needs Section \ref{an07cwiv1:precision}.\\
WIT \texttt{propagation}
& A solution of the defined leading system and the selected
 solution at the same \(\rho\); chart validity on \([a,b]\),
 \(M,N\le2\), fixed relative weights, logistic reactions;
 regular characteristics and finite distance. Any inherited
 probe/gap restrictions are those above; they are not used
 in the propagation estimate itself.
& L equations, not the original equations. E initial distance;
 O original continuation/chart validity. Fixed current weights
 are not an original-flow conclusion. Theorem
 \ref{an07cwiv1:propagation} handles weak first moments.\\
WIT \texttt{population\-transfer}
& \((\rho,\zeta)\in\mathcal P_*\), same selected half-mass clock;
 masses \(\le2\), chart supports throughout \(J\);
 \(\sup_J\dist\le\eps_*\), where
 \(\eps_*=\min(1/4,\sigma/4,p_*/(8C_1),\gamma_0/(8C_1))\),
 \(C_1=64R^2(1+4/g_*)\).
& L parameter rectangle, reference gap \(g_*/4\) and margins;
 E small distance only at \(s_n\); O original uniform comparison
 on \(J\). The weak-support replacement is
 Corollary \ref{an07cwiv1:witness}.\\
& Regular characteristics for the a.e. chain rule;
 selector-defined source rates, positive window lengths \(r/2\);
 \(n_\rho\ge1/4+\sigma/2\) and the strict source/total margins.
& L, WIT \texttt{generalrates,uniformtheorem}. No original
 pointwise rate conclusion at gate-hit atoms.\\
WIT \texttt{clockalignment}
& Leading logistic \(N\), \(0<N(\tau_a)<1\), same \(\lambda\);
 existence on the compared interval, and to \(\tau_h\) if it
 is called a hit. For a perturbed mass: absolute continuity,
 \(0<N<1\), \(N(0)=1/2\), and its logit defect.
& L for the auxiliary logistic; E centering in
 Proposition \ref{an07cwiv1:time}. O physical mass law/defect
 on later windows. No response monotonicity is inferred.\\
LP \texttt{diagonal-transfer}
& Compact offset set, any stated compact ratio set; selected
 limiting family with a common finite clearing time \(S\);
 actual initialized positive-noise family \(S_0(q)\);
 analytical physical centering.
& L clearing on the LPR parameter set; E initialization \(q^{10}\)
 and centering \eqref{an07cwiv1:center}. Uniformity is limited to
 sets on which these inputs are proved.\\
& Full original \(\sup_{\zeta,\rho,S\le s\le T}
 |q^{-2}(E-1/2)-\mathcal E|\to0\) for every finite \(T>S\);
 all background included, nonnegative original times.
& O compact-transfer. Nonnegative time follows from
 Proposition \ref{an07cwiv1:time}. A single witness window or
 static entry comparison cannot supply this premise.\\
LPR \texttt{combinedtheorem}
& Selected coherent leading branch, strong mass exactly one,
 own logistic half-mass clock, \(\mathcal P_*\);
 WIT uniform witness constants and LP's uniform clearing on
 its compact projection.
& L, already proved by LPR. Its selection/continuity/forward
 growth inputs are leading conclusions, not original ones.
 No initialized transfer or wider rectangle is inherited.\\
H2 A-error-transfer
& Same actual initialization and center; every fixed finite
 centered interval, uniform final compact ratio/offset region;
 full original output including background and membership flux.
& E entry and center; P bounds only to \(t_b\);
 O continuation and \(q^2\)-precision. Fixed ancestry avoids
 flux here, but future moving charts must account for it.\\
H2 A-rate-transfer
& Normalized source rates \(D_p^\tau/q^2\) on the fixed
 gate-separated learned \(J\); convergence or explicit errors
 below source and total margins, uniform parameters.
& O original window transfer. Section \ref{an07cwiv1:rates}
 gives an entry ledger, not this conclusion. Gate-hit atoms
 are outside the uniform rate scope.\\
AN7 items 1--6
& Reached-class original field and reactions; full residuals;
 \(q^2\)-precision; all long-delay terms; mass/response/amplitude
 phase and flux; every fixed error window and learned rate
 window; uniform parameters and explicit scale family.
& E entry/phase from WCS, with \(O(q^2\ell)=o(1)\);
 P RB residual/core/outside bounds; remaining reached-class
 continuation, output and source-rate precision are O.
 No canonical regime, power band or broad split reference.\\
\bottomrule
\end{longtable}\normalsize

\section{Static activity: current angles, signs and endpoints}
\label{an07cwiv1:activity}
\begin{proposition}[Exact classification]
\label{an07cwiv1:signs}
For \(\xi_q(\zeta)=(\sin q\zeta,-\cos q\zeta)\) and a
positive-mass label,
\[
 a_\theta\cdot\xi_q=\sin(q\zeta-\theta).
 \]
It is active exactly when
\(\theta\in(q\zeta-\pi,q\zeta)\pmod{2\pi}\). Both endpoints
have zero activation. For \(0\le\zeta\le R\), \(qR<\pi/2\),
every current Q2 label is inactive, as is every current Q1
label with \(\theta\ge q\zeta\). A weak-chart label is inactive
if \(q|u+\zeta|\le\pi/2\) (strict inequality gives a negative
input). At \(t_n\), all of \(\Gp\), including its endpoint labels,
are inactive for every fixed bounded \(\zeta\).
\end{proposition}
\begin{proof}
The identity proves the interval and its endpoints.
For a weak receiver the input is \(-\cos(q(u+\zeta))\).
WCS's reached chart has \(qL\to0\), proving the assertion for \(S\).
To include \(P_q\), use RCET's positive \(U_2\) and ratio,
restated with its late-time proof in SRR \texttt{tail}:
\[
 \frac{q|U_1|}{U_2}
 \le C(C_0e^{-b_*(t-t_0)/2}+q^2),\quad
 e^{-b_*(t_s-t_0)}\le Cq^{597/320}.
 \]
Thus at \(t_n\ge t_s\) this ratio is \(o(1)\), and
\[
 \frac{U\cdot\xi_q}{U_2}
 \le R\,\frac{q|U_1|}{U_2}-\cos(q\zeta)<0 .
 \]
All divisions use the positive-coordinate conclusion of RCET,
whose regularized endpoint statement includes \(b=\pi\).
\end{proof}
For negative offsets the unqualified Q2 statement is false:
\(\theta=\pi\) has input \(-\sin(q\zeta)>0\) when \(\zeta<0\).
The correct active Q2 slice is \((\pi+q\zeta,\pi]\).
This does not obstruct WIT, whose offsets are positive.
Statements about Q2 inactivity must refer to current angles,
or to a separately reached inactive chart.

For WIT's positive offsets the activity ledger is:
\begin{center}\small
\begin{tabular}{p{.26\textwidth}p{.65\textwidth}}\toprule
Fixed cohort & Exact test, not a mass-based sign inference\\\midrule
Core \(\KK\) & \(x<\zeta\); \(x=\zeta\) gives zero.
 At \(t_n\) WIT offsets have a fixed positive gap from
 \(X_*(n)\), so all core labels are active for small \(q\).\\
\(S,P_q\) & Inactive at \(t_n\), by Proposition
 \ref{an07cwiv1:signs}.\\
\(\Sigma\) & Recruited Q2 ancestry does not imply inactivity.
 A part in a strong chart is active when \(x<\zeta\).
 Later lower-side inputs are not excluded by an energy norm.\\
Strong complement \(\CC\) & Use its current angle in the exact
 active interval. No full-cohort inactivity conclusion.\\
Q3 interior, \(3\pi/2\) strip & Current Q3 inputs are active
 exactly above \(\pi+q\zeta\); current Q4 inputs are active.
 Transit/returned labels use the same interval test.\\
\(\pi\) strip & The part still in Q3 with current
 \(\theta>\pi+q\zeta\) is active; equality gives zero.
 Initial offset alone does not decide later activity.\\\bottomrule
\end{tabular}
\end{center}
The angle \(3\pi/2=-\pi/2\) has input \(\cos(q\zeta)>0\);
\(\pi\) has input \(-\sin(q\zeta)\), zero only at \(\zeta=0\).
At overlapping ancestry endpoints assign \(\pi\) to the primary
closed cohort and \(3\pi/2\) to the right-half endpoint if a
disjoint integral partition is needed. They have zero initial
label measure; their pointwise signs are nevertheless checked.
No gate selector is substituted for a positive ReLU value.

\section{The \texorpdfstring{\(q^{-2}\)}{q-inverse-squared} precision ledger at entry}
\label{an07cwiv1:precision}
Put \(z_\alpha=a_{\theta_\alpha}\cdot\xi_q\). For a fixed
cohort \(T\) define
\begin{equation}
 \mathcal Q_T(\zeta)=\int_Tm(z_\alpha)_+^2\dd\lambda_0,
 \qquad \mu_T=\int_Tm\dd\lambda_0 .
 \label{an07cwiv1:Qprobe}
\end{equation}
This probe moment is not the data second moment \(Q\).
RB \texttt{postalignment}, inherited on its completed interval,
gives \(|\theta-\psi|\le\delta_q=C_{\rm al}q(1+\ell)\)
for every label. Constants are fixed before \(q\).

\begin{lemma}[Exact output split and its decisive precision]
\label{an07cwiv1:split}
At \(t_n\), set \(T=\Sigma\cup\BB\), so \(f=f_\KK+f_T\). Exactly,
\begin{align}
 E_{\xi_q}-\tfrac12
 &=-\xi_q\cdot f+\tfrac12|f|^2, \notag\\
 E_{\xi_q}-E_{\xi_q}^{\KK}
 &=-\xi_q\cdot f_T+f_\KK\cdot f_T+\tfrac12|f_T|^2,\notag\\
 -\xi_q\cdot f_T&=-\mathcal Q_T+\mathcal A_T,\qquad
 |\mathcal A_T|\le\delta_q\sqrt{\mu_T\mathcal Q_T},\qquad
 |f_T|\le\sqrt{\mu_T\mathcal Q_T}.
 \label{an07cwiv1:exactsplit}
\end{align}
Here \(E^\KK=\tfrac12|\xi_q-f_\KK|^2\), not a
separately trained network. The core has \(|f_\KK|\le C_Rq\).
Since \(\mu_T(1+\ell)^2\to0\), uniformly over bounded offsets
\begin{equation}
 E_{\xi_q}-E_{\xi_q}^{\KK}=o(q^2)
 \quad\Longleftrightarrow\quad
 \mathcal Q_{\Sigma\cup\BB}=o(q^2).
 \label{an07cwiv1:equivalence}
\end{equation}
\end{lemma}
\begin{proof}
For one label the linear term equals
\(-m(z_\alpha)_+^2
 -m[(a_{\psi_\alpha}-a_{\theta_\alpha})\cdot\xi_q](z_\alpha)_+\).
Use \(|a_\psi-a_\theta|\le|\psi-\theta|\) and
Cauchy--Schwarz. The output bound follows the same way.
The core is on its fixed chart. S2CP and RB give
\(\mu_\Sigma\le C(q^5+q^{\Gamma_\nu}H+q^2K_Y^2)\),
\(\Gamma_\nu>3/10\), and \(\mu_\BB\le q^{11/10}\);
at entry \(H=O_n(1)\). Hence the mass-log product vanishes.
The error in replacing \(E-E^\KK\) by \(-\mathcal Q_T\)
is at most
\[
 C(q+\delta_q)\sqrt{\mu_T\mathcal Q_T}
                 +\tfrac12\mu_T\mathcal Q_T.
 \]
Young's inequality absorbs any fixed small multiple of
\(\mathcal Q_T\), with remainder \(Cq^2(1+\ell)^2\mu_T=o(q^2)\).
This proves both directions, including the uniform version.
\end{proof}
For aligned labels \(W\parallel U\) with the initialized positive
parallel orientation, the linear contribution is exactly
\(-\mathcal Q_T\le0\). This is not an absolute-error bound
or a sign assertion for the complete quadratic observable
or its source derivatives.

\paragraph{Core and inactive cohorts.}
Taylor expansion on the reached core chart, exactly as in
P \texttt{pa:probe-pop}, gives
\[
 (f_\KK)_1=qA_\KK+O_R(q^3),\quad
 (f_\KK)_2=q^2B_\KK+O_R(q^4),\quad
 E^\KK-\tfrac12=q^2\mathcal E_\KK+O_R(q^4).
 \]
This is a self-output calculation; it does not assume the
remaining network absent. The strong first-moment entry bound gives
\begin{equation}
 E^\KK(t_n)-\tfrac12
   =q^2\mathcal E_{\rho,\zeta}(s_n)+O_n(q^{2+\chi})+O(q^4).
 \label{an07cwiv1:corestatic}
\end{equation}
It pays core mass and radius errors. Inactive labels
contribute exactly zero, regardless of radial mass.
For \(\BB\) alone the crude cross terms are already paid:
\[
 q\mu_\BB=O(q^{21/10}),\quad
 \mu_\BB^2=O(q^{11/5}),\quad
 \delta_q\mu_\BB=O(q^{21/10}(1+\ell)).
 \]
The unpaid term is its squared active input.
For \(\Sigma\), whose raw mass has a weaker power, retain the
square-root expression in \eqref{an07cwiv1:exactsplit};
do not use these \(\BB\)-specific crude estimates.

\paragraph{Cohort ledger.}
Write \(L_q=1+\sqrt{\log\ell}\) for small \(q\). The secondary
exponent is exactly
\[
 \Gamma_\nu=1-5\lambda+\frac{2\lambda}{1-\lambda}
                   -\nu[10(1-\lambda)-2],\qquad
 \frac{289169}{924000}\le\Gamma_\nu
                    \le\frac{481}{1020}<\frac12 .
 \]
The lower bound is S2CP \texttt{gamma}; for the upper bound,
the unpenalized expression increases on the ratio box
because its derivative is \(-5+2/(1-\lambda)^2>0\), and its
value at \(49/100\) is \(481/1020\).
Every entry below is at \(t_n\); the mass estimates themselves
hold only on the passage. A negative displayed exponent describes
failure of an available upper bound to prove the target,
not a lower bound on the true contribution.
\small
\begin{longtable}{>{\raggedright\arraybackslash}p{.19\textwidth}>{\raggedright\arraybackslash}p{.39\textwidth}>{\raggedright\arraybackslash}p{.34\textwidth}}
\toprule Cohort & Proved size / exact contribution & \(q^2\)-precision status\\\midrule
\endhead
\(\KK\) & \eqref{an07cwiv1:corestatic}; radius \(q^\chi\),
 mass defect \(q^{11/10}\). & Paid, error \(O_n(q^{2+\chi})+O(q^4)\).\\
\(S,P_q\) & \(f=0\) exactly at entry. & Paid as direct carriers.
 Their training residuals still drive the core.\\
\(\Sigma\) & \(\mathcal Q_\Sigma\le
 C(E_\Sigma+q^2J_\Sigma)\le
 C(q^5+q^{\Gamma_\nu}H+q^4)\).
 & Open at fixed \(H\asymp n\): gives
 \(O_n(q^{\Gamma_\nu-2}+q^2)\) after division.
 Need \(\mathcal Q_\Sigma=o(q^2)\).\\
Strong complement \(\CC\) & \(\mathcal Q_\CC\le\mu_\CC
 \le Cq^{1+\Delta_\nu}\), \(\Delta_\nu>3/25\).
 Refined test below. & Mass-only bound gives \(Cq^{\Delta_\nu-1}\).
 Need \(\mathcal Q_\CC=o(q^2)\).\\
Q3 interior & \(\mathcal Q\le\mu\le Cq^{10}\)
 (GESG, inherited by RB). & Paid, \(O(q^8)\) after division.\\
\(\pi\) strip & \(\mathcal Q\le Cq^{1211/800}L_q+
 Cq^{11}\sqrt\ell\) (PISB/RB).
 & Inner part open: \(O(q^{-389/800}L_q)\).
 Need its \(\mathcal Q=o(q^2)\).\\
\(J_3\), \(3\pi/2\) inner zone & \(\mathcal Q\le
 Cq^{4769/4000}L_q\) (LSB/RB).
 & Open: \(O(q^{-3231/4000}L_q)\).
 Need its \(\mathcal Q=o(q^2)\).\\
Frozen middle/outer zones & \(\mathcal Q\le
 Cq^{11}\sqrt\ell\) (PISB, included in RB).
 & Paid, \(O(q^9\sqrt\ell)\). Do not count these masses twice
 inside the strip totals.\\\bottomrule
\end{longtable}\normalsize
The bound for \(\Sigma\) follows from
\(\sin(q\zeta-\theta)^2\le
C_R(\sin^2\theta+q^2\cos^2\theta)\).
Its currently upper-half active slice has \(0\le\theta<q\zeta\)
and thus \(\mathcal Q\le C_Rq^2\mu_\Sigma=o(q^2)\).
The exact remaining possible slice is
\begin{equation}
 \int_{\Sigma\cap\{\sin\theta<0\}}
       m[\sin(q\zeta-\theta)]_+^2\dd\lambda_0=o(q^2).
 \label{an07cwiv1:Sigmagap}
\end{equation}
SEP's entry chart at \(t_0\), or an unsigned distance/gain bound
after it, does not prove this slice stays empty. We neither
assume such a sign nor claim the actual slice has large mass.

\subsection{What the strong-complement tail actually pays}
To distinguish PROFILE's exponent from centered time \(s\), write
\[
 P=P_\rho=\Phi(\rho a),\quad a=\rho\Kfun(\rho a),\quad
 \alpha=\lambda P/2,\quad d=1/2-\alpha,\quad
 s_{\rm pr}=\frac{1-\lambda P}{1-\lambda},\quad p=3/s_{\rm pr}.
 \]
Here \(\Kfun(z)=\phi(z)+z\Phi(z)\).
S2CP's proved entry uses
\[
 \epsilon_h=\sqrt2 A_q^d(1-M_0)^\alpha
                  q^{10d-s_{\rm pr}},\qquad
 d_c=c_*q^{\chi+kP/2},\qquad
 I(t)=Q\int_{t_0}^t(1-c(v))\dd v.
 \]
\(M_0\) is MOMC's moving-overlap strong mass, not the present
core mass. The \(A_q\) and \(1-M_0\) factors have positive fixed
upper/lower bounds. The cutoff constant \(c_*>0\) is already fixed.
\begin{proposition}[A direct second-moment test of the existing tail]
\label{an07cwiv1:tailtest}
Use the proved S2CP entry law \(\eta_0=2m(t_0)\dd\lambda_0\),
\(\eta_0\{\xi>u\}\le C\min(1,(\epsilon_h/u)^p)\),
\(p=3/s_{\rm pr}>2(1+\nu)\), and cutoff \(d_c\).
The existing distance and capped gain imply, on the passage,
\begin{align}
 \mathcal Q_\CC
 &\le Cq^2\mu_\CC+
 C(\epsilon_h/d_c)^p
       \min\{1,Ad_c^2\}\max\{1,Ad_c^2\}^{1+\nu},
 \label{an07cwiv1:tailbound}\\
 A&=C_{\rm fixed}q^2 e^{\beta I},\quad \beta=\lambda/Q.
 \notag
\end{align}
The endpoint majorant \(A\asymp q^{2-k}\), \(k=10(1-\lambda)\),
gives the exponent
\begin{equation}
 \Xi_\nu=p(E-\chi)-R_g-\nu(R_g)_+,\qquad
 E=5(1-P)-s_{\rm pr},\quad R_g=k(1-P)-2-2\chi .
 \label{an07cwiv1:Xi}
\end{equation}
This test would pay the probe precision if \(\Xi_\nu>2\)
uniformly. In fact its exponent obeys \(\Xi_\nu\le915/604<2\)
at \(\lambda=49/100\) for every allowed \(P\); it cannot pay
the requested uniform precision.
\end{proposition}
\begin{proof}
Before physical exit, the distance envelope and core angle
\(O(q)\) give
\[
 (z_\alpha)_+^2\le Cq^2+C\min(1,A\xi^2).
 \]
After exit, \(A\xi^2\) has reached a fixed positive threshold;
\(I\) increases, so the same inequality holds after a fixed
enlargement depending on the physical exit threshold.
Multiply the second term by the gain
\(\max(1,A\xi^2)^{1+\nu}\) and integrate over \(\xi>d_c\).
For \(Ad_c^2\le1\), split at \(A^{-1/2}\).
The lower integral is at most
\(C A\epsilon_h^p d_c^{2-p}\), since \(p>2\);
the upper is \(C\epsilon_h^p A^{p/2}\), since \(p>2(1+\nu)\).
It is no larger than the first bound because
\((Ad_c^2)^{p/2-1}\le1\).
If \(Ad_c^2\ge1\), the \(2(1+\nu)\) tail moment gives
the bound directly. No \(4+2\nu\) moment is needed.

S2CP justifies the trajectory powers and \(Q\)-versus-\(\lambda\)
clock factors: \(\epsilon_h/d_c\asymp q^{E-\chi}\) and
\(Ad_c^2\asymp q^{-R_g}\). This proves \eqref{an07cwiv1:Xi}.
On the independent rectangle used by TAILREP, take the exact
corner \(\lambda=49/100,\ P=153/250\). Then
\[
 R_g=-303/2500,\qquad
 \Xi_\nu=1+\frac{11004659}{43757500}
        =\frac{54762159}{43757500}<2.
 \]
This value is independent of \(\nu\), since \(R_g<0\).
This is exact rational substitution, not a numerical sign test.
There is also a check on the actual endpoint parameter.
Put \(B=p(E-\chi)-1-R_g\), as in TAILREP. At fixed
\(\lambda=49/100\),
\[
 \partial_P B=\frac{1-\lambda}{20(1-P\lambda)^2}
 [200P^2\lambda^2-400P\lambda+297\lambda-100]<0 .
 \]
The bracket decreases for \(P\lambda<1\), and at \(P=1/2\)
it is \(50\lambda^2+97\lambda-100=-8093/200<0\).
Thus, for every \(P\in[1/2,153/250]\), including \(P_\rho\),
\[
 \Xi_\nu=1+B-\nu(R_g)_+
 \le1+B(49/100,1/2)
 =\frac{2967}{1510}-\frac9{20}=\frac{915}{604}<2.
 \]
The missing inequality fails for this particular available
majorant even on the endpoint of the actual model box.
This is not a lower bound for the original probe moment or
a counterexample to a sharper probe estimate.
\end{proof}
The bounded \(A_q\) factors and fixed \(n\) alter constants,
not this missing exponent margin. A sharper angular/gain
correlation is an S5c obligation, not proved here.

\subsection{Source-rate precision at the same entry state}
\label{an07cwiv1:rates}
Here use
\(R^{(p)}_\theta=\frac12\E_{P_p}[(X-f)X^T
\one_{\{a_\theta\cdot X>0\}}]\); the mixture half is included.
Define
\[
 g_p=a_\psi^TR^{(p)}a_\theta,\quad
 v_{\theta,p}=a_\psi^TR^{(p)}a_\theta^\perp,\quad
 v_{\psi,p}=(a_\psi^\perp)^TR^{(p)}a_\theta .
 \]
The exact source output derivative of a fixed cohort is
\begin{equation}
 \begin{split}
 \dot f_{T,p}=\int_Tm\{&
 2g_p a_\psi z_+
 +v_{\psi,p}a_\psi^\perp z_+\\
 &+v_{\theta,p}a_\psi\one_{\{z>0\}}
                       (a_\theta^\perp\cdot\xi_q)\}\dd\lambda_0,\qquad
 D_p^\tau=-(\xi_q-f)\cdot\sum_T\dot f_{T,p}.
 \end{split}
 \label{an07cwiv1:exactrate}
\end{equation}
Use the selector convention, or a.e. time derivatives; no
uniform pointwise assertion is made at gate-hit atoms.

\begin{lemma}[Outside rate charge under the passage bounds]
\label{an07cwiv1:ratecharge}
At entry, for \(p=1,2\) and any fixed \(T\),
\begin{align}
 |\xi_q\cdot\dot f_{T,p}|
 &\le C\{\mathcal Q_T+q^2(1+\ell)^2\mu_T\},\notag\\
 |\dot f_{T,p}|
 &\le C\{\sqrt{\mu_T\mathcal Q_T}+q(1+\ell)\mu_T\}.
 \label{an07cwiv1:ratechargeeq}
\end{align}
Thus the unpaid \(o(q^2)\) conditions in the static ledger
also suffice for outside source-rate charges at \(t_n\).
Q3 interior and frozen-zone charges are already \(o(q^2)\).
\end{lemma}
\begin{proof}
Loss bounds and the two cross-residual norms, applied to each
cluster separately by Cauchy--Schwarz, give bounded diagonal
entries and off-diagonal entries \(O(q)\) of \(R^{(p)}\),
uniformly in its receiver. Hence \(|g_p|\le C\), and RB
alignment gives
\[
 |v_{\theta,p}|+|v_{\psi,p}|
       \le C(|\sin\theta\cos\theta|+q+\delta_q).
 \]
On an active label \(|\sin\theta|\le z_++C_Rq\).
Also \(|a_\psi\cdot\xi_q|\le z_++\delta_q\).
Insert these estimates in \eqref{an07cwiv1:exactrate}.
The projected integrand is bounded by
\(C[z_+^2+(q+\delta_q)z_+
                         \delta_q(q+\delta_q)]m\);
the unprojected one by \(C[z_++q+\delta_q]m\).
Integrate, use Cauchy--Schwarz and then Young's inequality.
No receiver gate was differentiated. The scope is the passage.
\end{proof}

\begin{lemma}[Core source precision at entry]
\label{an07cwiv1:corerates}
On WIT's offset range, the core's self-output rate satisfies
\begin{equation}
 -(\xi_q-f_\KK)\cdot\dot f_{\KK,p}
       =q^2d_{p,*}(s_n)+O_n(q^{2+\chi})+O(q^4),
       \qquad p=1,2.
 \label{an07cwiv1:corerateeq}
\end{equation}
These velocities are driven by the full original training
field. Inactive weak carriers are not deleted from the residual.
\end{lemma}
\begin{proof}
We restate the sourcewise consequence of CFA's componentwise
proof, rather than infer it from a total-field error.
At a core receiver put
\(\mathcal S(x)=\int_\KK\Ffun(\rho x,\rho\widetilde x)\dd\mu\)
and \(Y=\int_\KK y\dd\mu\). The unnormalized equations are
\begin{align*}
 2x'_1&=-(1-M)x+e_{x1},\\
 2y'_1&=-(1-M)y-Y-K_O+e_{y1},\\
 2x'_2&=\rho\phi(\rho x)-\rho\mathcal S(x)
            +\lambda[V_O+(1-c)y]\Phi(\rho x)+e_{x2},\\
 2y'_2&=\rho(1-c)\Kfun(\rho x)+e_{y2}.
\end{align*}
CFA \texttt{P1stein,P1error,P2assembly,gaussian} yields these
cluster-by-cluster identities with the pointwise majorant in
its proof; using physical \(M\) removes the normalization
cost. Explicitly, at \(t_n\) the error terms are bounded by
\[
 C[q^2+J_G+q\sqrt{E_GJ_G}
 +\mu_\Gp\mathcal T
 +E_\Sigma+q^{-1}\sqrt{E_\Sigma J_\Sigma}
 +\mu_\BB/q] .
 \]
Here \(J_G\le q^2K_Y^2,\ E_G=O(1)\), SRR
\texttt{tail} gives superpolynomial \(\mathcal T\), S2CP gives
\(E_\Sigma=O(q^5+q^{\Gamma_\nu})\), and
\(\mu_\BB/q\le q^{1/10}\).
Since \(\Gamma_\nu>3/10\), the sum is \(O(q^{1/10})\).
Bounded \(V_O,K_O\) follows from SRR \texttt{weakmoments}
and WCS entry; this pays their \(q^2\) multiples.
This is a pointwise argument, not use of CFA's integrated allowance.

The source radial coefficients also satisfy
\[
 (\log m)'_1=1-M+O(q^2+\mu_\BB),\qquad
 (\log m)'_2=O(q^2+\mu_\BB).
 \]
Indeed the core target/donor \(P_1\) longitudinal moment is
\(1-M+O(q^2)\). The non-core \(P_1\) first output costs
\(J_G+q\sqrt{E_GJ_G}=O(q^2)\), and the second output is
multiplied by \(O(q)\), giving \(q^2K_G=O(q^2)\).
On \(P_2\), the core input is \(O(q)(|Z|+|X_2|)\),
the first residual is \(O(q)\) in \(L^2\), and the second
output coordinate is \(O(q)\); the contractions cost
\(O(q^2)\). An arbitrary outside output costs \(C\mu_\BB\)
by Gaussian moments. This uses exactly the mixture
normalization in \eqref{an07cwiv1:exactrate}.

SRR's moment comparison and WCS entry replace \(c,V_O,K_O\)
by \(n,nV_*(n),n\Kfun(U_*(n))\) with error \(O_n(q^\chi)\).
Core donor and receiver errors have the same order.
Thus each scaled core source velocity and reaction differs
from the coherent leading counterpart by \(O_n(q^\chi)\).
The core probe gap is fixed and positive at entry. Differentiate
its compact trigonometric expansion with these bounded source
velocities: the error is \(O(q^4)\), and the leading rate is
WIT \texttt{generalrates}. Its strong first-moment comparison
gives \eqref{an07cwiv1:corerateeq}.
\end{proof}
In the full rate retain the additional terms
\[
 f_T\cdot\dot f_{\KK,p}
 -(\xi_q-f_\KK)\cdot\dot f_{T,p}
 +f_T\cdot\dot f_{T,p}.
 \]
Since \(\dot f_{\KK,p}=O(q)\), Lemmas
\ref{an07cwiv1:split} and \ref{an07cwiv1:ratecharge} pay them
if \(\mathcal Q_T=o(q^2)\). The same unpaid cohorts remain.
An \(O(q^{1/10})\) training-field charge alone is not their
probe-rate proof. After \(t_b\), the passage bounds used in
Lemma \ref{an07cwiv1:ratecharge} require a new source.

\section{Leading-system extension with no weak support bound}
\label{an07cwiv1:leading}
This section concerns only P's defined leading population
system. Its positive measures have masses \(M,N\), reactions
\(1-M,\lambda(1-N)\), and fixed within-family relative weights.
Set
\[
 Y=\int y\dd\mu_s,\quad V=\int v\dd\mu_w,\quad
 K_s=\int\Kfun(\rho x)\dd\mu_s,\quad
 K_w=\int\Kfun(u)\dd\mu_w ,
 \]
\[
 \Kfun(u)=\phi(u)+u\Phi(u),\qquad
 \Ffun(t,r)=\phi(\min(t,r))+r\Phi(\min(t,r)).
 \]
The strong source equations are those in the proof of
Lemma \ref{an07cwiv1:corerates}, now exactly with
\((c,V_O,K_O)=(N,V,K_w)\) and zero remainders.
The weak equations are
\begin{align}
 2u'&=\phi(u)-\mathcal W(u)
       -[Y+(1-M)v]\Phi(u)-\lambda(1-N)u,\notag\\
 2v'&=\rho K_s-\lambda V-\lambda(1-N)v-(1-M)\Kfun(u),
 \quad \mathcal W(u)=\int\Ffun(u,\widetilde u)\dd\mu_w.
 \label{an07cwiv1:weaksystem}
\end{align}
Let \(D_s,D_w\) denote the two integral terms in WIT's distance
about the coherent atom \((X,Y_*),(U,V_*)\); write
\(m_s=|M-1|,\ m_w=|N-n_*|\) and
\(\dist=m_s+m_w+D_s+D_w\).

\begin{theorem}[Rate robustness with weak first moments]
\label{an07cwiv1:rate}
In WIT \texttt{ratebound}, replace compact weak support by
\[
 \int (|u|+|v|)\dd\mu_w<\infty .
 \]
Keep every other hypothesis: \(R\ge2,\ 0<\rho\le1\),
\(M,N\le2\), nonnegative measures, strong support and all
comparison coordinates in \([-R,R]^2\), \(0\le n_*\le1\),
\(|\zeta|\le R\), and \(\zeta-X\ge g_0>0\).
For WIT's selector-defined leading rates,
\begin{equation}
 |d_p-d_{p,*}|\le64R^2(1+g_0^{-1})\dist,\qquad
 |\mathcal E-\mathcal E_*|\le10R^2\dist.
 \label{an07cwiv1:ratebound}
\end{equation}
For regular characteristic solutions with locally bounded
first moments, WIT's a.e. chain rule remains valid.
\end{theorem}
\begin{proof}
The identities
\[
 0\le\partial_t\Ffun(t,r)=(r-t)_+\phi(t),\qquad
 |\partial_r\Ffun|\le1,\qquad
 0\le\Kfun(u)\le\phi(0)+|u|
 \]
show that all weak moments exist. Do not assert a uniform
bound for actual \(K_w,V\), which can be arbitrarily large.
Compare donors to the bounded reference:
\begin{gather*}
 |A-A_*|\le D_s+2Rm_s,\quad
 |B-B_*|\le2R^2\dist,\quad |Y-Y_*|\le D_s+Rm_s,\\
 |K_w-n_*\Kfun(U)|\le D_w+2Rm_w,\qquad
 |V-n_*V_*|\le D_w+Rm_w .
\end{gather*}
Here \(A=\int(\zeta-x)_+\dd\mu_s,\ B=\int y(\zeta-x)_+\dd\mu_s\),
\(A_*=\zeta-X\), and \(B_*=Y_*A_*\).
For a strong receiver put \(D=|x-X|+|y-Y_*|\).
Change donor moments at the actual receiver first, then move
the receiver in the bounded reference field. As in WIT,
without an actual weak supremum, this gives
\begin{equation}
 |x'_2-x'_{2,*}|\le2R(D+\dist),\qquad
 |y'_2-y'_{2,*}|\le\tfrac12D+R\dist.
 \label{an07cwiv1:strongdiff}
\end{equation}
The reference satisfies \(|x'_{2,*}|,|y'_{2,*}|\le R\);
for example
\(2x'_{2,*}=\lambda\Phi(\rho X)
[n_*V_*-X+(1-n_*)Y_*]\).
The reference receiver derivative of the overlap is bounded
on the strong chart. No bound on \(\partial_t\Ffun\) at an
unbounded weak receiver is used here.

Use WIT's exact formulas
\[
 d_1=-\tfrac12A(Y+K_w)+(1-M)\mathcal C,\quad
 |\mathcal C|\le28R^2,\qquad
 d_2=\int_{x<\zeta}(\zeta-A-y)x'_2\dd\mu_s
       +\int(\zeta-x)_+y'_2\dd\mu_s .
 \]
Split the lag product as
\(A[(Y+K_w)-(Y_*+n_*\Kfun(U))]
 +(A-A_*)[Y_*+n_*\Kfun(U)]\).
Use \(A\le4R\) and bounded reference lag. Its half costs
\(7R^2\dist\); the reaction term costs \(28R^2\dist\).

For the input part of \(d_2\), write \(c_x=\zeta-A-y\),
\(c_*=X-Y_*\). Then \(|c_x|\le6R,\ |c_*|\le2R\).
Integrating \(c_x(x'_2-x'_{2,*})\) costs at most
\(36R^2\dist\) by \eqref{an07cwiv1:strongdiff}.
The term \((c_x-c_*)x'_{2,*}\) costs at most \(4R^2\dist\).
The active-mass mismatch costs
\(2R^2(m_s+D_s/g_0)\), since
\(\mu_s\{x\ge\zeta\}\le D_s/g_0\).
The output part costs at most \((13/2)R^2\dist\):
use \eqref{an07cwiv1:strongdiff} under
\((\zeta-x)_+\), then compare \(A\) with \(A_*\).
The total is at most \((50+2/g_0)R^2\dist\), below
\eqref{an07cwiv1:ratebound}.
The energy bound is WIT's unchanged
\(4R|A-A_*|+|B-B_*|\le10R^2\dist\).

On a finite interval the strong velocities are integrably
bounded by the weak first moments. Differentiating the strong
positive parts is justified by domination. The level-set chain
rule sets the normal contribution to zero at almost every time
on a zero level, exactly as in WIT. No assumption that an atom
never hits a gate is needed.
\end{proof}

\begin{theorem}[First-moment finite-window propagation]
\label{an07cwiv1:propagation}
Let a regular solution of the defined leading system and the
selected coherent solution at the same \(\rho\) have masses
in \([0,2]\), strong support and all comparison coordinates
in \([-R,R]^2\), \(R\ge2\), on \([a,b]\). Assume only finite
initial weak first moments. Then
\begin{equation}
 \dist(s)\le e^{64R(s-a)}\dist(a),\qquad a\le s\le b.
 \label{an07cwiv1:propbound}
\end{equation}
Strong chart validity is presupposed; no original-flow
continuation assertion is part of this theorem.
\end{theorem}
\begin{proof}
The strong full velocity comparison, with the same donor-first
split as above, is
\[
 |x'-X'|+|y'-Y_*'|\le8R(D+\dist).
 \]
For weak receivers the corresponding absolute Lipschitz bound
is unnecessary and need not hold uniformly. Use a one-sided
scalar estimate. Define
\[
 A_t(u)=\tfrac12[\phi(u)-\mathcal W(u)-Y\Phi(u)
                                      -\lambda(1-N)u].
 \]
Since \(|Y|\le2R\), \(N\le2\) and the weak measure is positive,
\begin{equation}
 A_t'(u)=-\tfrac12[(u+Y)\phi(u)
            +\mathcal W'(u)+\lambda(1-N)]
       \le R+1\le2R\quad(u\in\mathbb R).
 \label{an07cwiv1:onesided}
\end{equation}
Here \(|u|\phi(u)\le1\) and
\(\mathcal W'(u)=\phi(u)\int(\widetilde u-u)_+\dd\mu_w\ge0\).
Retain the Gaussian factor; no negative-input restriction or
uniform weak receiver radius was inserted.
Comparing at bounded \(U\) costs at most \(3R\dist/2\).
The remaining input term satisfies
\[
 \tfrac12|1-M|\,|v|\Phi(u)
 \le\tfrac12|v-V_*|+\tfrac R2m_s ,
 \]
and the output term satisfies
\[
 \tfrac12|1-M|\Kfun(u)
 \le\tfrac12|u-U|+Rm_s .
 \]
Use the moment estimates in the output equation to obtain
\[
 D^+(|u-U|+|v-V_*|)
       \le8R(|u-U|+|v-V_*|+\dist).
 \]
This replaces WIT's weak absolute-velocity comparison and
proves exactly the distance inequality required.

Finite first moments stay locally finite on the stipulated
interval: \(|\Ffun(u,r)|\le\phi(0)+|r|\),
\(\Kfun(u)\le\phi(0)+|u|\), bounded strong moments and
\eqref{an07cwiv1:weaksystem} give a linear first-moment
Gronwall bound. This also justifies integration of the
characteristic inequalities by truncation and domination.
Reaction coefficients have absolute value at most one,
and relative weights are fixed in this leading system.
Consequently
\[
 D^+D_s\le(1+8R)D_s+16R\dist,\qquad
 D^+D_w\le(1+8R)D_w+16R\dist .
 \]
The two logistic mass differences cost at most three times
their sum on \([0,2]\). Thus
\(D^+\dist\le(1+40R)\dist\le64R\dist\).
Integrate. Zero family mass is covered by continuity.
\end{proof}

\begin{corollary}[WIT with strong-only compact support]
\label{an07cwiv1:witness}
WIT \texttt{populationtransfer} remains valid if compact weak
support is replaced by finite first weak moments.
Keep its leading-system, same-parameter and clock hypotheses,
strong support and comparison coordinates in \([-R,R]^2\),
\(M,N\le2\), regularity for a.e. derivatives, and
\(\sup_J\dist\le\eps_*\) with the unchanged \(\eps_*\).
Then all its mass, source-rate and total-rate inequalities
hold, and both rebound bounds are at least \(\gamma_0r/4\).
An entry distance \(\eps_*e^{-128Rr}\) at the left endpoint
of \(J\), with the stated strong chart validity, suffices.
\end{corollary}
\begin{proof}
Apply Theorem \ref{an07cwiv1:rate} with \(g_0=g_*/4\);
WIT's margins and integration argument are unchanged.
Theorem \ref{an07cwiv1:propagation} gives the last assertion.
These hypotheses define the leading-model class; none is
advertised as a discharged original-window hypothesis.
\end{proof}

\paragraph{How WCS enters this extension.}
Let \(d=|u-U_*(n)|+|v-V_*(n)|\). The fixed incoming law obeys
\[
 \eta_s\{d>1\}\le W_2^2=O_n(q^{1/10}),\qquad
 \int_{\{d>1\}}d\dd\eta_s\le W_2^2=O_n(q^{1/10}).
 \]
The second inequality uses \(d\le d^2\) there.
Current probability has the same estimates up to \(1+o(1)\);
the whole first moment is \(O_n(q^\chi)\).
This supplies entry without discarding the tail or declaring it
compact. A leading comparison initialized with the selected
core/weak state, optionally rescaled as in Section 1, has distance
\(O_n(e^{64R(s-s_n)}q^\chi)\) on each fixed interval where its
stipulated strong chart holds. This is a leading statement,
not an original-flow estimate on that interval.

\section{S5b/S5c handoff and fixed constant order}
\label{an07cwiv1:handoff}
Closed results are the clock conversion and interface,
the exact activity classification, the static/source-rate
precision reduction and paid cohorts, and the leading-model
first-moment extension. Complete original \(q^2\) probe
comparison remains open even at entry until the unpaid moments
in Section \ref{an07cwiv1:precision} are paid. No missing item
is added as a premise of an initialized theorem.

\small
\begin{longtable}{>{\raggedright\arraybackslash}p{.18\textwidth}>{\raggedright\arraybackslash}p{.75\textwidth}}
\toprule Consumer & Exact remaining deliverable\\\midrule
\endhead
S5b: state & On each fixed \([s_n,T]\), continue the original
 strong chart, physical masses, weak first/second moments and
 appropriate residual bounds. Prove a moment-stable
 original-to-leading field/reaction reduction with errors
 tending to zero and fixed \(e^{C(T-s_n)}\) losses.
 Weak supremum compactness is no longer a leading-consumer
 requirement. Growing-chart finite-noise errors and their
 moment integrals still require proof.\\
S5b: clock & Use \(\tau=\tau_h+s\) from
 \eqref{an07cwiv1:center}; prove \(N_S-n_\rho=o(1)\) on every
 compared interval, or its logit defect
 \(\int r_N/[N_S(1-N_S)]=o(1)\), with entry mismatch
 \(O_n(q^{1/8})\). Include flux if selected labels change.
 Do not center through the \(O(1)\) terminal passage-time error.\\
S5b: weak activity & Keep \(q|u+\zeta|<\pi/2\) on the desired
 offsets, or pay exact active weak output and rate.
 Finite first moments in the leading model alone do not
 prove finite-noise weak inactivity.\\
S5c: output & For every fixed centered interval and final
 compact offset/ratio region, prove the full remainder
 \(-\xi_q\cdot f_T+f_\KK\cdot f_T+|f_T|^2/2=o(q^2)\),
 uniformly. A sufficient form under continued alignment is
 \(\sup\mathcal Q_{\Sigma\cup\BB}=o(q^2)\) and
 \(\sup\mu_T(1+\ell)^2=o(1)\).
 Pay separately the strong complement, possible lower-side
 secondary slice, \(\pi\) inner strip and \(J_3\), including
 transit labels. Passage mass exponents alone do not pay them.\\
S5c / rate stage & On WIT's gate-separated \(J\), prove
 \(\xi_q\cdot\dot f_{T,p}=o(q^2)\) and the cross terms in
 \eqref{an07cwiv1:exactrate} are \(o(q^2)\).
 Lemma \ref{an07cwiv1:ratecharge} suffices if its matrix and
 alignment bounds are continued; otherwise supply direct
 projected source bounds. Retain a sourcewise core
 velocity/reaction reduction, not just a total-field error.\\
WIT / H2 & Obtain \(\sup_J\dist\le\eps_*\), original normalized
 rate errors below fixed margins (or \(o(1)\) after division
 by \(q^2\)), and arbitrary fixed-window error convergence.
 Rate claims at gate-hit atoms remain excluded.\\
LP / LPR & Use the same physical center
 \(t_h^{\rm phys}=\tau_h/\mu_1^2\). Supply LP's every-finite-\(T\)
 compact-transfer premise after common clearing \(S\);
 then its diagonal conclusion applies, without a quantitative
 horizon rate. LPR is already a leading theorem on
 \(\mathcal P_*\), not original-flow transfer.\\\bottomrule
\end{longtable}\normalsize

Keep WPT/RB's constants, then WBR/MP/SRR/WCS's constants and
fixed \(n\), unchanged. Next fix the final compact parameter/
offset region (use \(\mathcal P_*\) for the common witness),
a finite centered interval, and \(R\ge2\) containing the
strong chart and reference coordinates with slack. Fix the
positive reference probe gap on \(J\), then
\(64R^2(1+g_0^{-1})\), \(\eps_*\), and
\(e^{64R(T-s_n)}\); choose \(q\) last.
New leading constants depend only on \(R,g_0\) and already
fixed witness margins, not residual guards. Original entry
constants may depend on fixed \(n\) and earlier guards; they
multiply displayed vanishing powers. No earlier chart or guard
is enlarged after seeing \(q\).

\paragraph{Boundary of this increment.}
The false unrestricted-offset Q2 inactivity proposal was
corrected to the exact sign range; WIT's positive offsets are
unaffected. The tail test records a precision gap in the
available implication, not a failure of the initialized theory.
No original continuation past \(t_b\), new strip mass theorem,
gate-hit rate transfer, AN8 assembly, S6 result or
canonical-regime claim is made.
\end{document}
